\subsection{Macaulay modules and flat algebras}

PROPOSITION 9. — Let \(\rho: A \to B\) be a homomorphism of Noetherian rings, M a finitely generated A-module and N a finitely generated B-module, flat over A. Denote by \({}^{a}\rho: \operatorname{Spec}(B) \longrightarrow \operatorname{Spec}(A)\) the map associated with \(\rho\). The following conditions are equivalent:

\begin{enumerate}
\def\labelenumi{(\roman{enumi})}
\item
  the B-module M ⊗A N is Macaulay;
\item
  the \((\kappa(\mathfrak{p})\otimes_{\mathrm{A}}\mathrm{B})\) -module \(\kappa(\mathfrak{p})\otimes_{\mathrm{A}}\mathrm{N}\) is Macaulay for every \(\mathfrak{p}\in\mathrm{Supp}_{\mathrm{A}}(\mathrm{M})\) , and the \(\mathrm{A}_{\mathfrak{p}}\) -module \(\mathrm{M}_{\mathfrak{p}}\) is Macaulay for every \(\mathfrak{p}\in{}^{a}\rho(\mathrm{Supp}_{\mathrm{B}}(\mathrm{N}))\) ;
\item
  for every maximal ideal of \(\mathrm{Supp}_{\mathrm{B}}(\mathrm{N})\) whose inverse image \(\mathfrak{p}\) in A belongs to \(\mathrm{Supp}_{\mathrm{A}}(\mathrm{M})\) , the \(\mathrm{A}_{\mathfrak{p}}\) -module \(\mathrm{M}_{\mathfrak{p}}\) and the \((\kappa(\mathfrak{p}) \otimes_{\mathrm{A}} \mathrm{B})\) -module \(\kappa(\mathfrak{p}) \otimes_{\mathrm{A}} \mathrm{N}\) are Macaulay.
\end{enumerate}

If moreover the B-module N is faithfully flat, these conditions imply that the A-module M is Macaulay.

Let q be a prime ideal of B belonging to the support of \(M \otimes_{A} N\) . Set \(\mathfrak{p} = \rho^{-1}(\mathfrak{q})\) . Since the module \((\mathrm{M} \otimes_{\mathrm{A}} \mathrm{N})_{\mathfrak{q}}\) is identified with \(M_{p} \otimes_{A_{p}} N_{q}\) , the modules \(M_{p}\) and \(N_{q}\) are nonzero, and the same is true of \(\kappa(\mathfrak{p}) \otimes_{A} N_{\mathfrak{q}}\) ( \(n^{\circ} 6\) , remark). The \(A_{p}\) -module \(N_{q}\) , being isomorphic to a module of fractions of \(N_{p}\) , is flat and one has the equalities

\[
\operatorname{prof} _ {\mathrm{B} _ {\mathfrak {q}}} \left(\left(\mathrm{M} \otimes_ {\mathrm{A}} \mathrm{N}\right) _ {\mathfrak {q}}\right) = \operatorname{prof} _ {A_ {\mathfrak {p}}} \left(\mathrm{M} _ {\mathfrak {p}}\right) + \operatorname{prof} _ {\mathrm{B} _ {\mathfrak {q}}} (\kappa (\mathfrak {p}) \otimes_ {\mathrm{A}} \mathrm{N} _ {\mathfrak {q}})
\]

\[
\dim_ {B _ {q}} \left(\left(M \otimes_ {A} N\right) _ {q}\right) = \dim_ {A _ {p}} \left(M _ {p}\right) + \dim_ {B _ {q}} \left(\kappa (p) \otimes_ {A} N _ {q}\right)
\]

(§ 1, n° 6, prop. 11, b) and remark), in which each term is an integer ≥ 0. Taking into account the fact that the B \(_{q}\) -module κ(p) ⊗A N \(_{q}\) is Macaulay if and only if it is Macaulay as a (κ(p) ⊗A B \(_{q}\) )-module (n° 1, example 4), one deduces the equivalence of the following two conditions:

(α) the Bq-module (M ⊗A N)q is Macaulay;

(β) the \(A_{p}\) -module \(M_{p}\) and the \((\kappa(\mathfrak{p})\otimes_{A}B_{\mathfrak{q}})\) -module \(\kappa(\mathfrak{p})\otimes_{A}N_{\mathfrak{q}}\) are Macaulay. Let us now prove that (iii) implies (i). Let n be a maximal ideal of B belonging to the support of \(M\otimes_{A}N\) , and \(\mathfrak{p}=\rho^{-1}(\mathfrak{n})\) . From what precedes, we have \(\mathfrak{p}\in\operatorname{Supp}_{A}(M)\cap{}^{a}\rho(\operatorname{Supp}_{B}(N))\) ; condition (iii) and the remark of no. 6 imply that condition (β) above is satisfied with q=n. It follows that the \(B_{n}\) -module \((M\otimes_{A}N)_{n}\) is Macaulay, whence (i).

The implication (ii)⇒ (iii) is clear; let us prove that (i) implies (ii). Suppose the B-module M ⊗\_A N is Macaulay. Let p be an element of Supp\_A(M). We may assume that the (κ(p) ⊗\_A B)-module κ(p) ⊗\_A N is nonzero, that is, there exists a prime ideal q of Supp\_B(N) lying over p (no. 6, remark). The B\_q-module (M ⊗\_A N)\_q is Macaulay (no. 1, example 3); it follows from the implication (α) ⇒ (β) proved previously and from the remark of no. 6 that the A\_p-module M\_p and the (κ(p) ⊗\_A B)-module (κ(p) ⊗\_A N) are Macaulay, whence (ii).

If moreover N is faithfully flat over A, we have \(\kappa(\mathfrak{p})\otimes_{\mathrm{A}}\mathrm{N}\neq 0\) for every \(\mathfrak{p}\in\operatorname{Spec}(\mathrm{A})\) , whence \({}^{a}\rho(\operatorname{Supp}_{\mathrm{B}}(\mathrm{N}))=\operatorname{Spec}(\mathrm{A})\) (no. 6, remark), so that (ii) implies that M is Macaulay.

COROLLARY 1.--- Let \(\rho: A \to B\) be a local homomorphism of Noetherian local rings, M a nonzero finitely generated A-module and N a nonzero finitely generated B-module which is a flat A-module. For the B-module M \(\otimes_{A}N\) to be Macaulay, it is necessary and sufficient that the A-module M be Macaulay and that the B/ \(\mathfrak{m}_{A}\)B-module N/ \(\mathfrak{m}_{A}\)N be Macaulay.

Indeed, N is a faithfully flat A-module since \(N/m_{A}N\) is nonzero (I, § 3, no. 1, definition 1).

COROLLARY 2.-- Let \(\rho: A \to B\) be a homomorphism of Noetherian rings making \(B\) a faithfully flat \(A\)-module and let \(M\) be an \(A\)-module of finite type. For the \(B\)-module \(M \otimes_{A} B\) to be Macaulay, it is necessary and sufficient that the \(A\)-module \(M\) be Macaulay and that \(\kappa(\mathfrak{p}) \otimes_{A} B\) be a Macaulay ring for every \(\mathfrak{p} \in \operatorname{Supp}(M)\) .

COROLLARY 3.--- Let A be a Noetherian ring, B a finite flat A-algebra, and M a finitely generated Macaulay A-module. The B-module M⊗ \(_{A}\) B is Macaulay.

Indeed, for every \(\mathfrak{p} \in \operatorname{Spec}(A)\) , the ring \(\kappa(\mathfrak{p}) \otimes_{A} B\) is a finite \(\kappa(\mathfrak{p})\)-algebra, hence is Macaulay (no. 5, example 1), and we apply prop. 9.

COROLLARY 4.--- Let A be a Noetherian ring, J an ideal of A, and M a finitely generated A-module. Denote by \(\hat{A}\) and \(\hat{M}\) the separated completions of A and M for the J-adic topology, and S the multiplicative subset \(1 + J\) from A. Consider the following conditions:

\begin{enumerate}
\def\labelenumi{(\roman{enumi})}
\item
  the A-module M is Macaulay
\item
  the \(\widehat{\mathrm{A}}\) -module \(\widehat{\mathrm{M}}\) is Macaulay;
\item
  the \(\mathrm{S}^{-1}\mathrm{A}\) -module \(\mathrm{S}^{-1}\mathrm{M}\) is Macaulay;
\item
  for every maximal ideal \(\mathfrak{m} \in \operatorname{Supp}(\mathrm{M}) \cap \mathrm{V}(\mathrm{J})\) , the \(\mathrm{A}_{\mathfrak{m}}\) -module \(\mathrm{M}_{\mathfrak{m}}\) is Macaulay;
\item
  for every prime ideal \(\mathfrak{p} \in \operatorname{Supp}(M)\) such that \(\mathfrak{p} + J \neq A\) , the \(A_{\mathfrak{p}}\) -module \(M_{\mathfrak{p}}\) is Macaulay and the ring \(\kappa(\mathfrak{p}) \otimes_{A} \widehat{A}\) is Macaulay.
\end{enumerate}

Conditions (ii) through (v) are equivalent, and are implied by (i). When the ideal J is contained in the radical of A, conditions (i) through (v) are equivalent.

We know that (i) implies (iii) (no. 1, example 3), and (iii) is identical to (i) when J is contained in the radical of A (since the elements of S are then invertible).

The ring \(\widehat{\mathrm{A}}\) is Noetherian (III, § 3, no. 4, prop. 8); it is identified with the completion of S \(^{-1}\) A for the topology S \(^{-1}\) J-adic, and the \(\widehat{\mathrm{A}}\) -module \(\widehat{\mathrm{M}}\) is identified with the completion of S \(^{-1}\) M for the topology S \(^{-1}\) J-adic (III, § 3, no. 5, prop. 12). Consequently, to prove the equivalence of conditions (ii) through (v), one may replace \(A\) by S \(^{-1}\) A, J by S \(^{-1}\) J and M by S \(^{-1}\) M; in other words, one may assume that J is contained in the radical of A. The A-module \(\widehat{\mathrm{A}}\) is then faithfully flat (loc. cit., prop. 9).

It is clear that (v) implies (iv) and that (iv) implies (i).

(i)⇒(ii): Let m be a maximal ideal of A; then m is a maximal ideal of  above m, and every maximal ideal of  is obtained in this way (III, § 3, no. 4, prop. 8). The ring κ(m) ⊗A  is a field, hence a Macaulay ring; if the A-module M is Macaulay, then so is the Â-module M by the implication (iii)⇒(i) of prop. 9.

(ii)⇒(v): The \(\hat{A}\) -module \(\hat{M}\) is isomorphic to \(M\otimes_{A}\hat{A}\) (III, § 3, no. 4, th. 3); if it is Macaulay, it follows from prop. 9, (i)⇒(ii) that \(\kappa(\mathfrak{p})\otimes_{A}\hat{A}\) is a Macaulay ring for every \(\mathfrak{p}\in\mathrm{Supp}(M)\) and that the A-module M is Macaulay.

PROPOSITION 10. - Let \(\rho: A \to B\) be a homomorphism of Noetherian rings making B a flat A-module. The following conditions are equivalent:

\begin{enumerate}
\def\labelenumi{(\roman{enumi})}
\item
  B is a Macaulay ring;
\item
  for every prime ideal q of B, the rings \(A_{\rho^{-1}(\mathfrak{q})}\) and \(\kappa(\rho^{-1}(\mathfrak{q}))\otimes_{A}B\) are Macaulay;
\item
  for every maximal ideal n of B, the rings \(A_{\rho^{-1}(\mathfrak{n})}\) and \(\kappa(\rho^{-1}(\mathfrak{n})) \otimes_{A} B\) are Macaulay.
\end{enumerate}

If moreover B is faithfully flat over A, these conditions imply that A is a Macaulay ring.

This is the special case M = A, N = B of Prop. 9.

COROLLARY 1.--- Every algebra that is finite and flat over a Macaulay ring is a Macaulay ring.

COROLLARY 2. Let A be a Macaulay ring and n a positive integer; then \(\mathrm{A}[\mathrm{X}_1,\dots ,\mathrm{X}_n]\) and \(\mathrm{A}[[\mathrm{X}_1,\dots ,\mathrm{X}_n]]\) are Macaulay rings.

It suffices to treat the case \(n = 1\). The ring A{[}T{]} is Noetherian (A, VIII, § 1, n° 4, cor. 1), and, for every field \(k\) , the ring \(k[T]\) is a Macaulay ring (no. 5, example 2); consequently, the ring A{[}T{]} is Macaulay (prop. 10) and the ring A{[}{[}T{]}{]} is Macaulay (cor. 4 of prop. 9).

COROLLARY 3.--- Every finitely generated algebra over a Macaulay Noetherian ring is catenary.

Indeed, such an algebra is a quotient of a polynomial ring over a Macaulay ring, hence a quotient of a Macaulay ring (cor. 2), and consequently is catenary (no. 2).

